\documentclass[12pt,UTF8,fontset=windows]{ctexart}

\usepackage[a4paper,margin=1in]{geometry}
\usepackage{amsmath}
\usepackage{amssymb}
\usepackage{booktabs}
\usepackage{array}
\usepackage{tikz}
\usetikzlibrary{arrows.meta,positioning}
\usepackage{hyperref}
\usepackage{setspace}
\usepackage{enumitem}
\usepackage{caption}
\hypersetup{colorlinks=true,linkcolor=black,citecolor=black,urlcolor=blue}

% 约1.2倍行距
\setstretch{1.2}
\setlength{\parskip}{2pt plus 1pt}
\setlength{\parskip}{0pt plus 1pt}
\setlength{\parindent}{2em}
\setlength{\abovecaptionskip}{4pt}
\setlength{\belowcaptionskip}{2pt}
\setlength{\intextsep}{4pt}
\setlength{\textfloatsep}{8pt}
\setlength{\floatsep}{6pt}
\setlength{\abovedisplayskip}{8pt}
\setlength{\belowdisplayskip}{8pt}
\setlength{\abovedisplayshortskip}{4pt}
\setlength{\belowdisplayshortskip}{4pt}
\setlength{\jot}{4pt}

% OOF公式行距略宽
\everydisplay{\setstretch{1.2}}
\setlength{\emergencystretch}{3em}
\sloppy
\setlist{nosep,leftmargin=2em,labelindent=0pt}

% 章标题四号(14pt)加粗居左；小标题小四(12pt)加粗居左
\ctexset{
  section={
    beforeskip=18pt, afterskip=10pt,
    format={\Large\bfseries\raggedright}, numbering=true,
  },
  subsection={
    beforeskip=12pt, afterskip=6pt,
    format={\normalsize\bfseries\raggedright}, numbering=true,
  },
}

\begin{document}

% ═══ 首页 ═══
\begingroup
\centering
{\Large\bfseries Home Credit Open Solution 的受控复现}\\[6pt]
{\large\bfseries 性能增益来自特征表示、动态行为，还是模型集成？}\\[10pt]
{\normalsize 工业数据分析课程项目 \hspace{2em} \today}
\par
\endgroup
\vspace{4pt}

% ═══ 1  问题背景与受控复现设计 ═══
\section{问题背景与受控复现设计}

\subsection{任务与动机}

Home Credit Default Risk 竞赛要求预测贷款申请人是否会出现还款困难。每名申请人由
\texttt{SK\_ID\_CURR} 唯一标识，二元标签 \texttt{TARGET} 在出现还款困难时取 1，评价指标为 ROC
曲线下面积（AUC）——衡量模型是否给高风险申请人赋予更高预测概率的排序能力。

这份数据不是一张扁平的单表，而是一组关系型表格。中心 application 表每名申请人占一行，含
基本信息、收入、资产、家庭等静态字段；其余五张表分别记录：\texttt{bureau}（征信局历史，含
贷款金额、状态、逾期金额）、\texttt{previous\_application}（历史贷款申请，含申请金额、审批
结果）、\texttt{installments\_payments}（分期还款明细，含应还/实还金额、逾期天数）、
\texttt{POS\_CASH\_balance}（POS 与现金贷款月度余额，含逾期状态 SK\_DPD）、
\texttt{credit\_card\_balance}（信用卡月度余额，含额度、使用金额）。核心建模难点不只是"训
练一个分类器”，更在于如何把变长的历史记录转化为申请人级别的特征——这一步“表示构造”
决定了模型最终能"看到"哪些信息。

开源方案 \texttt{minerva-ml/open-solution-home-credit} 经历了六个有名称的版本（表~\ref{tab:versions}）：
Chestnut（单表基线）、Seedling（群组聚合与多模型）、Blossom（多表历史聚合）、Tulip（群体相对
与近期行为）、Sunflower（清洗修复与动态特征）、Four Leaf Clover（OOF stacking），公开榜 AUC
从约 0.74 升至约 0.80。但这些版本不能直接当作受控实验来比较——相邻版本间特征工程、
清洗规则、模型类别与超参数往往同时变化，分数提升无法干净归因。

\begin{table}[h]
\centering\small
\caption{作为项目参照的历史 Open Solution 版本。}\label{tab:versions}
\begin{tabular}{p{0.15\textwidth}p{0.38\textwidth}p{0.35\textwidth}}
\toprule
版本 & 主要变化 & 在本项目中的角色 \\
\midrule
Chestnut & application 单表基线 & S1 基线 \\
Seedling & application 群组聚合与更多模型 & S2 群组聚合 \\
Blossom & 多张历史表的聚合特征 & S3 历史信息 \\
Tulip & 群体相对特征与近期行为 & S4 更精细的表示 \\
Sunflower & 清洗修复与动态特征 & B2/S5 清洗与动态 \\
Four Leaf Clover & 在 OOF 预测上做 stacking & S6 模型集成 \\
\bottomrule
\end{tabular}
\end{table}

本项目的核心问题：当开源解答分数提升时，究竟是哪一类新增的信息结构带来了增益？ 三个研究
问题（RQ）为：

\begin{enumerate}
\item RQ1：从 Chestnut 到 Blossom 的早期增益来自哪里？候选：群组聚合、业务比例、多表历史。
\item RQ2：从 Blossom 到 Sunflower 的后期增益来自哪里？候选：相对位置、最近窗口、清洗修复、动态趋势。
\item RQ3：充分的特征工程之后，stacking 能否继续提升？
\end{enumerate}

\subsection{实验矩阵}

\begin{table}[h]
\centering\small
\caption{受控特征阶段与对应原版本。}\label{tab:matrix}
\begin{tabular}{llll}
\toprule
阶段 & 对应版本 & 特征内容 & 目的 \\
\midrule
S1   & Chestnut       & application 数值+类别（78 列） & 基线 \\
S2   & Seedling       & S1 + 群组聚合原值（379 列）    & 通用群组聚合 \\
B1   & 桥接           & S1 + 业务比例+EXT\_SOURCE 汇总（85 列） & 业务特征桥接 \\
S3   & Blossom        & B1 + 五表历史聚合（126 列）     & 历史信息 \\
S4   & Tulip          & S3 + 群体相对位置+最近窗口（184 列） & 更精细表示 \\
B2   & 桥接           & S4，聚合前清洗（184 列）        & 清洗顺序 \\
S5   & Sunflower      & B2 + 短期/长期比值+趋势（191 列）& 动态行为 \\
S6   & Four Leaf Clover & 在一级 OOF 预测上做 stacking & 预测互补 \\
\bottomrule
\end{tabular}
\end{table}

\subsection{评价协议与折安全}

所有 LightGBM 阶段使用相同五折分层划分（\texttt{data/folds.csv}，seed=2026）、相同
LightGBM 参数（lr=0.03，num\_leaves=31，feature\_fraction=0.20，reg\_lambda=100，
n\_estimators=3000，early\_stopping=100）。主指标为折外 AUC（OOF AUC）。增益判定采用
描述性判据：平均 $\Delta>0$ 且至少 4/5 折同向，称为"稳定增益"；否则为"不稳定/证据不足"。

每条 OOF 预测来自一个没见过该申请人的模型：
\[\text{OOF}=\{(\texttt{SK\_ID\_CURR}_i, y_i, \hat p_i, \texttt{fold\_id}_i)\}_{i=1}^{n}.\]
相比训练集 AUC，OOF AUC 更可靠——每条预测都来自一个没有见过该申请人的子模型。

折安全是核心不变量：每个群组聚合都只在当前训练折内拟合，对验证折或
测试集的申请人仅作为查找表使用。若验证折出现训练折未见过的 group，则用训练折上该聚合的
全局均值填补。构造群组聚合时绝不使用标签。任何与折相关的统计均遵循"先在训练折
\texttt{fit}、再分别 \texttt{transform}"模式——验证行永远不会参与生成用于预测自身的统计量。
这一要求使实验成本上升（需要折内循环中的特征构造），但每一组对比的可信度都建立在这一基础
之上——没有它，"稳定增益"的判据就失去了意义。

\subsection{可复现产物与 S1/S2 详情}

每个阶段均输出四个核心产物：\texttt{oof.parquet}（含 SK\_ID\_CURR、TARGET、fold\_id 与
prediction）、\texttt{fold\_metrics.csv}、\texttt{feature\_names.txt}（特征列名，须与模型
输入列完全一致）、\texttt{config.yaml}；加 \texttt{--predict-test} 时额外输出
\texttt{submission.csv}。实现校验每名训练申请人恰好出现一次、预测值有限且在 $[0,1]$ 内、
折标识与 \texttt{folds.csv} 一致。

S1 仅使用 application 表的 78 个特征，若干已知异常编码转为缺失
（\texttt{DAYS\_EMPLOYED=365243}、\texttt{CODE\_GENDER=XNA} 等），刻意排除比例特征、群组
聚合、历史表特征与动态窗口，作为参照点。S2 在 S1 之上加入 fold-safe 的群组聚合原值（301 个
特征，379 列总计），聚合涵盖贷款金额、年金、收入、商品价格、外部评分、年龄相关天数、家庭
人数与车龄等核心变量，使用 min、mean、max、sum、var、median 函数，主要分组键包括
"学历$\times$性别""家庭状况$\times$学历""家庭状况$\times$性别"及若干职业与地区组合。
S2 只保留群组聚合原值，不构造 $x-\overline{x}_{g(i)}$ 形式的差值特征——差值保留到 S4。

为了检验 Seedling 式的特征空间是否特别依赖非线性树模型，我们在同一套 379 维 S2 特征上额外
训练了一个带 $L_2$ 正则的 Logistic 回归。数值变量在训练折内做中位数填补与标准化；类别变量
做 one-hot 编码，转换时忽略未见类别。该模型不是 S2 的主结果，而是一个桥接实验，同时也是
后续 S6 stacking 阶段可能的输入之一。诊断性 S2-full（含群组差值特征，801 列）OOF AUC 相对
S1 提升 $+0.0027$，表明信号在相对位置而非绝对统计量。

表~\ref{tab:official} 汇总所有已提交阶段的 Kaggle 官方 public/private AUC。官方榜单整体与
OOF 方向一致：B1 和 S3 在 public/private 上均明显高于 S1/S2 系列，S4/S5 均高于 S3。但
public/private 的采样噪声和排行榜划分会让相邻阶段出现轻微换序（如 S4 的 public 最高而 S5
的 private 最高）。S6 两个 stacking 版本已完成官方提交，结果同样低于 S5 单模型。因此本文把官方分数作为外部有效性检查，归因以固定 folds 下 OOF 配对对比
为主。

\begin{table}[h]
\centering\small
\caption{全阶段官方榜单汇总。}\label{tab:official}
\resizebox{\textwidth}{!}{%
\begin{tabular}{lrrrr}
\toprule
阶段 & 特征数 & OOF AUC & Public AUC & Private AUC \\
\midrule
S1 & 78  & 0.757646 & 0.74770 & 0.74418 \\
S2 & 379 & 0.757405 & 0.74656 & 0.74471 \\
S2-full & 801 & 0.760329 & 0.74920 & 0.74859 \\
S2-LR & 379 & 0.743054 & 0.73544 & 0.73049 \\
B1 & 85  & 0.768159 & 0.76716 & 0.76307 \\
S3 & 126 & 0.785109 & 0.78900 & 0.78562 \\
S4 & 184 & 0.786169 & 0.79294 & 0.78666 \\
B2 & 184 & 0.786196 & 0.79164 & 0.78686 \\
S5 & 191 & 0.786342 & 0.79137 & 0.78696 \\
S6-avg & 4 & 0.777088 & 0.77651 & 0.77079 \\
S6-stack & 4 & 0.780383 & 0.78104 & 0.77552 \\
\bottomrule
\end{tabular}
}
\end{table}

% ════════════════════════════════════════════════════════════════════
% 第 2 章  RQ1：早期增益归因
% ════════════════════════════════════════════════════════════════════
\section{RQ1：从基础特征到关系型表示}

\subsection{对比设计与结果}

RQ1 将早期（Chestnut $\to$ Blossom）拆为四组受控对比：S2$-$S1（群组聚合原值）、B1$-$S1
（业务比例+EXT\_SOURCE）、S3$-$B1（五表历史聚合）、S2-LGBM$-$S2-Logistic（树非线性）。

B1 是桥接实验：仅新增 7 个特征——三个业务比例
$\frac{\text{AMT\_ANNUITY}}{\text{AMT\_INCOME\_TOTAL}}$、
$\frac{\text{AMT\_CREDIT}}{\text{AMT\_INCOME\_TOTAL}}$、
$\frac{\text{AMT\_CREDIT}}{\text{AMT\_ANNUITY}}$，
以及 EXT\_SOURCE\_1/2/3 的 mean、min、max、std——把"业务比例"从"多表历史"中剥离。
S3 在 B1 之上接入五表历史聚合（41 个新特征）。

\begin{table}[h]
\centering\small
\caption{RQ1 各阶段结果。}\label{tab:rq1-stages}
\begin{tabular}{llrrrr}
\toprule
阶段 & 模型 & 特征数 & OOF AUC & 折 AUC 标准差 & 同向折数 \\
\midrule
S1 & LightGBM & 78  & 0.757646 & 0.002460 & --- \\
S2 & LightGBM & 379 & 0.757405 & 0.002358 & 1/5 \\
B1 & LightGBM & 85  & 0.768159 & 0.001864 & 5/5 \\
S3 & LightGBM & 126 & 0.785109 & 0.001543 & 5/5 \\
S2-Logistic & Logistic & 379 & 0.743054 & 0.003636 & --- \\
\bottomrule
\end{tabular}
\end{table}

\begin{table}[h]
\centering\small
\caption{RQ1 受控配对对比。}\label{tab:rq1-deltas}
\begin{tabular}{l p{0.32\textwidth} r c p{0.14\textwidth}}
\toprule
对比 & 含义 & $\Delta$OOF AUC & 同向折数 & 稳定性 \\
\midrule
S2$-$S1 & 群组聚合原值 & $-0.000241$ & 1/5 & 不稳定 \\
B1$-$S1 & 业务比例+EXT\_SOURCE & $+0.010513$ & 5/5 & 稳定 \\
S3$-$B1 & 五表历史聚合 & $+0.016950$ & 5/5 & 稳定 \\
S2-LGBM$-$S2-Logistic & 树非线性 & $+0.014350$ & 5/5 & 稳定 \\
\bottomrule
\end{tabular}
\end{table}

\begin{figure}[h]
\centering
\begin{tikzpicture}[
  box/.style={draw,rounded corners,align=center,inner sep=5pt,font=\footnotesize},
  arr/.style={-{Stealth[length=2.4mm]},thick}
]
\node[box,fill=gray!12]   (s1) at (0,0)      {S1\\0.7576};
\node[box,fill=blue!8]    (s2) at (4.5,1.5)  {S2\\0.7574};
\node[box,fill=yellow!18] (b1) at (4.5,-1.5) {B1\\0.7682};
\node[box,fill=green!14]  (s3) at (9.0,-1.5) {S3\\0.7851};
\draw[arr] (s1) -- (s2) node[midway,above,sloped,font=\footnotesize]{$\Delta=-0.0002$（不稳）};
\draw[arr] (s1) -- (b1) node[midway,below,sloped,font=\footnotesize]{$\Delta=+0.0105$（稳）};
\draw[arr] (b1) -- (s3) node[midway,above,font=\footnotesize]{$\Delta=+0.0170$（稳）};
\end{tikzpicture}
\caption{RQ1 OOF AUC 增益树。}\label{fig:gain-tree}
\end{figure}

\subsection{增益归因分析}

群组聚合原值为何无效。 S2 的 $\Delta=-0.000241$，五折仅 1 折正向。301 个群组均值
（如"学历$\times$性别群体的收入均值"）与底层个体变量高度共线，且对群体内所有人取相同值——LightGBM
本就能直接在底层变量上做分裂，群组均值几乎无新增信息。在 \texttt{feature\_fraction=0.20} 的列采样下，
大量低信息特征反而稀释了每棵树采到有效特征的概率。诊断性的 S2-full（加入群组差值特征）OOF AUC 提升
$+0.0027$，表明信号在"个体偏离群体均值的程度"而非"群体绝对统计量"——这正是把差值特征留到 S4 的原因。

业务比例的高效率。 B1 仅新增 7 个特征即获得 $+0.01051$，单位特征增益 $1.50\times10^{-3}$，
为 RQ1 最高（表~\ref{tab:rq1-density}）。EXT\_SOURCE 是外部征信机构已"预消化"的风险评分，
其 mean/min/max 捕捉多来源共识，std 捕捉分歧——几个外部评分相互矛盾本身就是风险信号。三个业务
比例把"能否负担这笔贷款"显式表达：年金/收入、贷款/收入刻画偿债压力，贷款/年金刻画还款期限结构。

多表历史聚合是最大跃升。 S3$-$B1 的 $+0.01695$ 是整个复现链中绝对增益最大的一步。
历史表提供直接的行为证据：bureau 的在外负债规模、installments 的逾期比例刻画还款纪律、
previous\_application 的批准率反映被金融机构接受的程度。把变长历史记录聚合到申请人粒度
是本题最关键的表示步骤——这是"把关系型历史转化为申请人级表示"这一动作的回报。

树非线性的独立贡献。 S2-LGBM$-$S2-Logistic 的 $+0.01435$ 衡量完全相同特征空间下模型
非线性带来的价值。Logistic（OOF AUC 0.743）甚至低于 S1 LightGBM（0.758），说明"好特征+线性模型"
不一定胜过"树模型+朴素特征"——特征表示与模型能力是两个正交维度。

\begin{table}[h]
\centering\small
\caption{各特征块信息密度（$\Delta$AUC $\div$ 新增特征数）。}\label{tab:rq1-density}
\begin{tabular}{lrrr}
\toprule
特征块 & 新增特征数 & $\Delta$OOF AUC & 每特征增益 \\
\midrule
S2 群组聚合 & 301 & $-0.000241$ & $\approx 0$ \\
B1 业务比例 & 7   & $+0.010513$ & $+1.50\times10^{-3}$ \\
S3 历史聚合 & 41  & $+0.016950$ & $+4.13\times10^{-4}$ \\
\bottomrule
\end{tabular}
\end{table}

跨折标准差从 S1 的 0.002460 降至 B1 的 0.001864 再降至 S3 的 0.001543——有价值特征
不仅抬升均值，还压低五折间波动。B1、S3 均为五折全部同向。

\subsection{RQ1 结论与原始方案对照}

早期增益主要来自多表历史聚合（S3$-$B1，$+0.01695$，五折全同向），其次为业务比例与
EXT\_SOURCE（B1$-$S1，$+0.01051$）。普通群组聚合原值单独无贡献（S2$-$S1 $\approx 0$）。
"多做 groupby 就会更好"在本受控设置下不成立，群组信息真正的价值需以相对位置形式（S4）引入。

把上述结果放回原始版本历史，可以看到受控复现的价值。原始 Seedling 公开榜提升约 $+0.005$，
但那是"群组聚合 $+$ 群组差值 $+$ 更多模型"一并上线的结果。受控 S2 单独保留群组聚合原值
得增益约零——Seedling 被笼统归功于"groupby 特征"，受控对比表明真正起作用的是同期引入的
相对差值与模型扩展。同时，"好特征+线性模型"不一定胜过"树模型+朴素特征"——在 379 维上
Logistic（0.743）低于 S1 LightGBM（0.758），说明特征表示与模型能力是两个正交维度，本报告
的主线刻意固定模型、只改特征，正是为了把特征表示的贡献单独量化出来。

% ════════════════════════════════════════════════════════════════════
% 第 3 章  RQ2：后期增益归因
% ════════════════════════════════════════════════════════════════════
\section{RQ2：后期增益来自更精细的信息组织？}

\subsection{对比设计}

RQ2 将 Blossom $\to$ Sunflower 的复合升级拆为三组独立对比：S4$-$S3（群体相对位置+最近窗口）、
B2$-$S4（聚合前清洗修复）、S5$-$B2（短期/长期比值+趋势）。B2 是本报告最重要的桥接实验：
与 S4 使用完全相同的 184 列特征，唯一区别是历史表聚合在清洗后的原表上计算。

\subsection{结果与归因}

\begin{table}[h]
\centering\small
\caption{RQ2 各阶段结果与配对对比。}\label{tab:rq2}
\begin{tabular}{lrrrc}
\toprule
阶段 & 特征数 & OOF AUC & $\Delta$ vs 父阶段 & 同向折数 \\
\midrule
S3（Blossom）            & 126 & 0.785109 & ---          & --- \\
S4（Tulip，群体+最近）   & 184 & 0.786169 & $+0.001060$  & 4/5 \\
B2（清洗修复，桥接）     & 184 & 0.786196 & $\approx 0$  & 2/5 \\
S5（Sunflower，动态特征）& 191 & 0.786342 & $\approx 0$  & 3/5 \\
\bottomrule
\end{tabular}
\end{table}

\begin{figure}[h]
\centering
\begin{tikzpicture}[
  box/.style={draw,rounded corners,align=center,inner sep=5pt,font=\footnotesize},
  arr/.style={-{Stealth[length=2.4mm]},thick}
]
\node[box,fill=gray!12]   (s3) at (0,0)     {S3\\0.7851};
\node[box,fill=blue!8]    (s4) at (4.5,0)   {S4\\0.7862};
\node[box,fill=yellow!18] (b2) at (9.0,1.2) {B2\\0.7862};
\node[box,fill=green!14]  (s5) at (9.0,-1.2){S5\\0.7863};
\draw[arr] (s3) -- (s4) node[midway,above,font=\footnotesize]{$\Delta=+0.0011$（稳）};
\draw[arr] (s4) -- (b2) node[midway,above,sloped,font=\footnotesize]{$\approx 0$（不稳）};
\draw[arr] (b2) -- (s5) node[midway,left,font=\footnotesize]{$\approx 0$（不稳）};
\end{tikzpicture}
\caption{RQ2 OOF AUC 增益树。}\label{fig:rq2-gain-tree}
\end{figure}

群体相对位置与最近窗口的稳定贡献。 S4 新增 58 个特征：30 个群体相对位置（个人值减去同
职业/同学历/同性别群体均值及绝对差值）和 28 个最近窗口行为（最近 1/3/5 笔贷款申请、最近
10/50 条分期还款和 POS 消费记录）。$\Delta=+0.00106$，四折同向，稳定但量级仅为 RQ1 历史聚合的
1/16。机制上，"同一年收入 5 万元，对洗碗工和医生意味着完全不同的还款能力"——相对位置将信息从
绝对值转化为在同类人中的位置。类似地，最近窗口特征将信息从"全量均值"转化为"最近
趋势"：一个申请人在 installments 上过去 5 年整体逾期率不高，但最近 10 笔全部逾期——全量
均值会掩盖这一恶化的信号，而最近窗口能将其暴露出来。S4 中群体相对位置只保留差值
（$x-\bar{x}_g$）和绝对差值（$|x-\bar{x}_g|$），不保留群组均值原值——因为 RQ1 已证明群组
原值无增量信息，把原值与差值分开避免了把后者的功劳错记到前者头上。

清洗修复为何无效。 B2$-$S4 的 $\Delta\approx 0$（2/5 同向），并非实现错误，而是
清洗目标列不在聚合范围内。B2 清洗 bureau 的 DAYS\_* 列、previous 的 DAYS\_* 列、
credit\_card 的 AMT\_DRAWINGS\_* 列——但 S3/S4 的聚合只需要 AMT\_CREDIT\_SUM/DEBT/OVERDUE、
DAYS\_CREDIT、AMT\_APPLICATION/CREDIT、DAYS\_DECISION、CNT\_PAYMENT、AMT\_BALANCE、
AMT\_CREDIT\_LIMIT\_ACTUAL、SK\_DPD。所有被清洗列均不在聚合范围内。这在方法论上是有效的
实验结果：Sunflower 的清洗修复，其价值依赖于聚合列覆盖清洗目标列。

动态比值与趋势无稳定增益。 S5 新增 7 个特征（4 个短期/长期比值+3 个 POS SK\_DPD 趋势），
$\Delta=+0.00015$，仅 3/5 折同向。原因有三：比值与窗口特征高度冗余（LightGBM 可在相邻分裂中
隐式构造）；SK\_DPD 序列 97\% 为零，对稀疏信号拟合线性趋势斜率方差极大；本实验的特征数量远少于
原方案 Sunflower（数十列窗口比值和趋势），弱信号的累积效应无法体现。此外，原方案 Sunflower 在数十个列上计算窗口比值和趋势，而
非本实验仅限定的 4 个比值 $+$ 3 个趋势——可能依靠大量弱信号的累积效应来获得可测量提升，
这也是原方案 CV 值更高的一个技术解释。

\subsection{与原始方案的对照}

原始 Tulip 宣称 CV $+0.0065$，我们的受控 S4 仅 $+0.0011$——差额源自原方案同期加入的更丰富聚合
与手工特征（特征从约 200 增至数百），我们的 S4 严格控制仅打开两块（+58 列）。原始 Sunflower
的 $+0.0045$ 增益，根据 B2 和 S5 结果，几乎不能归因于清洗修复或动态比值/趋势——真正增益来源
可能在于同期加入的大规模特征扩展（聚合列扩至几乎所有数值列），而非清洁修复或动态分析本身。

\subsection{RQ2 结论}

后期增益中唯一可确认的来源是群体相对位置和最近窗口（S4$-$S3，$+0.0011$，四折同向）。
聚合前清洗在本实验的聚合列覆盖范围内无增益。短期/长期比值与趋势特征无稳定增益。
后期精细化收益仅为早期关系型表示增益的约 1/15——当核心关系型信息已被提取后，进一步精细化
的实际收益高度依赖特征覆盖广度。原始 Tulip 宣称的 CV $+0.0065$ 中仅约 $+0.0011$ 可归因于
相对位置与近期行为这一特定的信息组织方式，其余来自同期加入的更丰富聚合与手工特征。原始
Sunflower 的 $+0.0045$ 增益，根据 B2 和 S5 结果，几乎不能归因于清洗修复或动态特征——真正
来源可能在于大规模特征扩展（聚合列扩至几乎所有数值列）。

\[
+0.0169_{\text{历史聚合}} \;>\; +0.0105_{\text{业务比例}} \;\gg\; +0.0011_{\text{群体+最近}} \;\approx\; 0_{\text{清洗}} \;\approx\; 0_{\text{动态}}.
\]

% ════════════════════════════════════════════════════════════════════
% 第 4 章  RQ3：Stacking 能否继续提升
% ════════════════════════════════════════════════════════════════════
\section{RQ3：特征纵向叠加后，Stacking 能否继续提升？}

\subsection{实验设计}

S6 的输入是四个一级模型的 OOF 预测矩阵
$Z = [p_{\text{lr}}^{\text{s2}}, p_{\text{lgb}}^{\text{s3}}, p_{\text{lgb}}^{\text{s4}},
p_{\text{lgb}}^{\text{s5}}] \in \mathbb{R}^{n\times4}$（$n=307{,}511$）。

简单平均： $\hat p_i^{\text{avg}} = \frac{1}{4}\sum_{j=1}^4 Z_{i,j}$，无需任何训练。

二层 CV 下的 L2-Logistic Stacking： 复用 \texttt{folds.csv} 五折划分，对每折用其余
四折的 $(Z,y)$ 拟合 $L_2$-Logistic 元模型（penalty=l2, solver=newton-cholesky,
max\_iter=1000, class\_weight=balanced），对该折做预测，五折拼接得到 stacking OOF。

\subsection{结果与归因}

\begin{table}[h]
\centering\small
\caption{S5 与两种 stacking 方法的 OOF AUC 和官方榜单。}\label{tab:rq3-stages}
\setlength{\tabcolsep}{3pt}
\begin{tabular}{lrrrrr}
\toprule
方法 & OOF AUC & Public & Private & $\Delta$ vs S5 & 同向折数 \\
\midrule
S5（最优单模型）          & 0.786342 & 0.79137 & 0.78696 & ---           & --- \\
简单平均                  & 0.777088 & 0.77651 & 0.77079 & $-0.009254$   & 0/5 \\
L2-Logistic stacking      & 0.780383 & 0.78104 & 0.77552 & $-0.005959$   & 0/5 \\
\bottomrule
\end{tabular}
\end{table}

两项对比五折全部为负——stacking 不仅无增益，反而显著退步；官方 public/private 也均低于 S5。

\begin{table}[h]
\centering\small
\caption{一级模型 OOF 预测相关矩阵。}\label{tab:rq3-corr}
\begin{tabular}{lrrrr}
\toprule
& S2-LR & S3    & S4    & S5 \\
\midrule
S2-LR & 1.000 & 0.726 & 0.720 & 0.721 \\
S3    & 0.726 & 1.000 & 0.979 & 0.978 \\
S4    & 0.720 & 0.979 & 1.000 & 0.991 \\
S5    & 0.721 & 0.978 & 0.991 & 1.000 \\
\bottomrule
\end{tabular}
\end{table}

一级预测高度相关是退步的根本原因。 S3/S4/S5 三者 $\rho=0.978\text{--}0.991$，输出预测
几乎相同——三个 LightGBM 几乎在犯相同的错误，提供的互补信息极少。S2-LR 与 LightGBM 系列适度
互补（$\rho\approx0.72$），但其自身 AUC 仅为 0.743，比 S5 低 0.043。元模型
$\hat p=\sigma(\beta_0+\sum\beta_j p_j)$ 在 $p_3,p_4,p_5$ 高度共线时，$\beta_2,\beta_3,\beta_4$
的估计方差极大，退化成一个对 $p_5$ 和 $p_{\mathrm{lr}}$ 的加权平衡器——而 $p_{\mathrm{lr}}$ 信号
质量太差，拖低了任何包含它的加权。

\begin{table}[h]
\centering\small
\caption{Logistic stacking 元模型系数。}\label{tab:meta-coeffs}
\begin{tabular}{lr}
\toprule
输入 & 系数 \\
\midrule
$p_{\text{lr}}^{\text{s2}}$ (S2-LR) & $+1.716$ \\
$p_{\text{lgb}}^{\text{s3}}$ (S3)   & $+1.800$ \\
$p_{\text{lgb}}^{\text{s4}}$ (S4)   & $+2.290$ \\
$p_{\text{lgb}}^{\text{s5}}$ (S5)   & $+3.516$ \\
截距 & $-1.709$ \\
\bottomrule
\end{tabular}
\end{table}

元系数排序符合直觉：S5 最高权重（3.52），S2-LR 最低（1.72）。Logistic stacking 确实识别出
了 S5 是最优单模型。但二层 CV 的严格性导致在每折子集上，S2-LR 偶尔表现微弱互补价值，使元模型
不忍将其权重压至零；在全集 OOF 上评估时，这种微弱优势不足以抵消 S2-LR 引入的噪声。此外 sigmoid
的非线性映射和输入尺度差异也使加权组合不严格等价于直接使用 S5。

\subsection{与原始方案的对照}

原始 Four Leaf Clover 报告约 $+0.0025$ CV 增益（0.7950 $\to$ 0.7975），其 stacking 输入
更多样——在多个不同特征子集上的 LightGBM（如一个只用 bureau、一个只用 application），即
"横向切分"，预测相关性显著低于本实验 S3$\to$S4$\to$S5 的纵向叠加。当一级模型预测夹角
足够大时，stacking 能取长补短；当它们几乎共线时则不能。此外原方案 CV 值更高（0.795 vs
本实验的 0.786），更高的基底意味着更多特征列和更大的模型多样性。

本实验的受控 stacking 明确告诉我们：在特征纵向叠加范式下，即使把从 S2 到 S5 所有级别的模型
预测都交给 stacking，也无法从高度同质的 LightGBM 系列中提取额外增益。stacking 本身不是药方
——其价值高度依赖一级预测之间的互补性，而非一级预测的数量或特征数。

\subsection{RQ3 结论}

简单平均和 L2-Logistic stacking 均未能超越最优单模型 S5（退步 $-0.006$ 至 $-0.009$，
五折全部同向为负）。退步原因：S3/S4/S5 高度相关（$\rho>0.978$），几无互补信号；S2-LR 虽适度
互补但自身 AUC 过低。Stacking 不保证增益——其价值取决于一级预测的互补性，而非一级预测数量。

% ════════════════════════════════════════════════════════════════════
% 第 5 章  总结与结论
% ════════════════════════════════════════════════════════════════════
\section{总结与结论}

\subsection{完整证据链}

\begin{table}[h]
\centering\small
\caption{全部受控阶段概览。}\label{tab:full-chain}
\begin{tabular}{lllrr}
\toprule
阶段 & 新增内容 & 特征数 & OOF AUC & $\Delta$ vs 父阶段 \\
\midrule
S1   & application 基准特征 & 78   & 0.757646 & --- \\
S2   & +群组聚合原值 & 379  & 0.757405 & $-0.000241$ \\
B1   & +业务比例+EXT\_SOURCE & 85 & 0.768159 & $+0.010513$ \\
S3   & +五表历史聚合 & 126  & 0.785109 & $+0.016950$ \\
S4   & +群体相对位置+最近窗口 & 184 & 0.786169 & $+0.001060$ \\
B2   & +聚合前清洗（无效）& 184 & 0.786196 & $\approx 0$ \\
S5   & +短期/长期比值+趋势 & 191 & 0.786342 & $\approx 0$ \\
S6-avg   & 简单平均 stacking & 4 & 0.777088 & $-0.009254$ \\
S6-stack & L2-Logistic stacking & 4 & 0.780383 & $-0.005959$ \\
\bottomrule
\end{tabular}
\end{table}

\begin{figure}[h]
\centering
\begin{tikzpicture}[
  box/.style={draw,rounded corners,align=center,inner sep=4pt,font=\scriptsize},
  arr/.style={-{Stealth[length=2.0mm]},thick}
]
\node[box,fill=gray!12]    (s1)    at (0,0)          {S1\\0.7576};
\node[box,fill=yellow!18]  (b1)    at (3.2,2.6)      {B1\\0.7682};
\node[box,fill=green!14]   (s3)    at (6.5,1.8)      {S3\\0.7851};
\node[box,fill=blue!8]     (s4)    at (9.8,3.2)      {S4\\0.7862};
\node[box,fill=purple!12]  (b2)    at (9.8,1.2)      {B2\\0.7862};
\node[box,fill=orange!20]  (s5)    at (13.2,2.2)     {S5\\0.7863};
\node[box,fill=red!12]     (avg)   at (13.2,3.8)     {S6-avg\\0.7771};
\node[box,fill=red!12]     (stack) at (13.2,0.6)     {S6-stk\\0.7804};
\draw[arr] (s1) -- (b1) node[midway,above,sloped,font=\scriptsize]{$+0.0105$稳};
\draw[arr] (b1) -- (s3) node[midway,above,font=\scriptsize]{$+0.0169$稳};
\draw[arr] (s3) -- (s4) node[midway,above,sloped,font=\scriptsize]{$+0.0011$稳};
\draw[arr] (s4) -- (b2) node[midway,left,font=\scriptsize]{$\approx 0$};
\draw[arr] (b2) -- (s5) node[midway,above,font=\scriptsize]{$\approx 0$};
\draw[arr,red!70!black] (s5) -- (avg)  node[midway,right,font=\scriptsize]{$-0.0093$};
\draw[arr,red!70!black] (s5) -- (stack) node[midway,right,font=\scriptsize]{$-0.0060$};
\end{tikzpicture}
\caption{从 S1 到 S6 的完整证据链。黑色箭头为稳定增益，红色箭头为负效应。}\label{fig:final-chain}
\end{figure}

\subsection{对三个 RQ 的最终回答}

RQ1——早期增益来自哪里？ 从 S1（0.758）到 S3（0.785）的 $+0.027$ 中：最大
稳定来源是多表历史聚合（$+0.01695$），其次是业务比例与 EXT\_SOURCE（$+0.01051$）；
普通群组聚合原值无贡献（$\approx 0$）。

RQ2——后期增益来自哪里？ 从 S3（0.785）到 S5（0.786）的 $+0.0012$ 微升中：唯一
确认的来源是群体相对位置和最近窗口（$+0.00106$）。清洗修复与本实验聚合范围不匹配，无增益。
动态比值与趋势无稳定增益。

RQ3——Stacking 能否提升？ 简单平均和 Logistic stacking 均未能超越 S5（退步
$-0.006$ 至 $-0.009$），因为 S3/S4/S5 预测高度相关（$\rho>0.978$）。在特征纵向叠加范式下，
stacking 不产生增益。

\subsection{工作假设验证}

"最大增益来自关系型与时序型特征表示"——成立。 多表历史聚合与业务比例贡献了可解释增益
的约 96\%。"群组聚合弱于历史表特征"——成立。 普通群组聚合原值增益为零，S4 相对位置
也仅为历史聚合的 1/15。"stacking 仅带来较小边际增益"——更保守：不仅无增益，反而退步。

\subsection{方法论反思}

受控复现的方法论本身就是一项重要产出。桥接实验 B1 和 B2 把原本纠缠的改动拆开：B1 将"业务
比例"从"多表历史"中干净剥离，避免了把后者的功劳重复计算到前者头上；B2 揭示了原方案文档未
标注的事实——清洗修复在特定聚合列选择下不产生增益。折安全的严格执行使每一组对比的可信度
建立在杜绝信息泄漏的基座上。"均值正 $+$ 至少 4/5 折同向"的判据在 $n=5$ 折下避免了 $p$ 值
不可靠的问题。受控复现天然只能回答"在限定特征集下某机制产生多大增益"，每个增量都是该特征块
在受控条件下的下界估计——当特征空间更丰富时某些机制增益可能更大，但本文不对此做出断言。此外，
受控复现与"原样复刻"有本质差异：前者追求干净的因果推断，后者追求分数的完全匹配。本文
选择前者，因此每个增量都是该特征块在受控条件下的下界估计。

\subsection{结论}

本文在固定五折划分、固定 LightGBM 超参数、严格折安全的前提下，依次打开 8 个特征块，将整套
Open Solution 从 OOF AUC 0.758 到 0.786 的提升分解为可归因的配对对比。

实验结果表明：（1）最大性能提升来自关系型特征表示——将五张变长历史表聚合为申请人级
行为统计，贡献了 $+0.01695$ 的稳定增益，是整条链中信息增量最大的一步。（2）业务语义特征
效率极高——仅 7 个特征即贡献 $+0.01051$，单位特征信息密度约为历史聚合的 3.6 倍。
（3）后期精细化收益远小于早期表示构造——群体相对位置和最近窗口的稳定增益仅为
$+0.00106$，清洗修复与动态特征无法被确认，合计后期增益仅为早期的约 1/15。
（4）在特征纵向叠加范式下，stacking 不产生增益——简单平均与 Logistic stacking 均低于
最优单模型 S5，因 LightGBM 系列 OOF 预测高度相关（$\rho>0.978$），互补信号不足。

从 0.74 到约 0.80 的整段旅程中，大约 60--70\% 可归因于"把关系型历史转化为申请人级统计表示"
这一核心问题——这是表示学习问题，而非模型调参问题。业务语义特征贡献了约 20\%，剩余 10--15\%
分散在群体相对位置、最近行为等更细颗粒度的工程优化中。本实验最终以 S5（191 维，OOF AUC
0.786342，Kaggle private 0.78696）为最优单模型，完成了从 S1 到 S6 的全部受控对比。

从 0.74 到约 0.80 的整段旅程中，大约 60--70\% 可归因于"把关系型历史转化为申请人级统计
表示"这一核心问题——这是表示学习问题，而非模型调参问题。业务语义特征贡献了约 20\%，
剩余 10--15\% 分散在群体相对位置、最近行为等更细颗粒度的工程优化中，这些优化在原方案更大
的特征空间下可能有更显著的效果，但在本实验的最小化受控设定下其独立贡献有限。本文也提醒：
受控复现天然只能回答"在限定特征集下某机制产生多大增益"——每个增量都是该特征块在下界
估计，当特征空间更丰富时某些机制（尤其是 S4 和 S5）的增益可能更大。
\end{document}
